% !PRINTED 323
\weberpage{323}
%<<CONT
ecke ändern ihr Vorzeichen, so oft der Punkt $x$ durch eine Seite
des Coordinatendreiecks hindurchgeht.

Mit Hülfe dieser nicht homogenen Relation kann die Einheit
linear und homogen durch $x_1, x_2, x_3$ dargestellt und dadurch
jede nicht homogene Function in eine homogene verwandelt
werden.

Die Coordinaten $x_1, x_2, x_3$ sind nach diesen Bestimmungen
zunächst nicht unabhängige Variable, sondern an die erwähnte
Relation gebunden. Wir erweitern ihre Bedeutung aber jetzt so,
dass wir sie als unabhängige Variable betrachten können, wenn
wir festsetzen, dass wir sie nur in homogenen Formen oder
Gleichungen benutzen wollen;
und dass nicht $x_1, x_2, x_3$ selbst,
sondern nur die Verhältnisse $x_1 : x_2 : x_3$ die Lage des Punktes $x$
bestimmen sollen. Dann können wir von der erwähnten Relation
zwischen den $x_1, x_2, x_3$ ganz absehen, und die $x_1, x_2, x_3$ als unabhängige Variable betrachten, für die nur die einzige Werthcombination $x_1 = 0, x_2 = 0, x_3 = 0$ ausgeschlossen ist. Die
Werthe $h x_1, h x_2, h x_3$ bestimmen dann für jedes von Null verschiedene $h$ einen und denselben Punkt $x$. Genau gesagt, sind
dann die Coordinaten $x_1, x_2, x_3$ des Punktes $x$ nicht die oben
beschriebenen Abstände von den Coordinatenaxen selbst, sondern irgend welche Grössen, die mit diesen Abständen proportional sind.

Jede lineare Substitution

\weq{1}{\begin{aligned}
x_1 &= \alpha_1^{(1)} y_1 + \alpha_1^{(2)} y_2 + \alpha_1^{(3)} y_3 \\
x_2 &= \alpha_2^{(1)} y_1 + \alpha_2^{(2)} y_2 + \alpha_2^{(3)} y_3 \\
x_3 &= \alpha_3^{(1)} y_1 + \alpha_3^{(2)} y_2 + \alpha_3^{(3)} y_3,
\end{aligned}}

\noindent deren Determinante

\weq{2}{r = \Sigma \pm \alpha_1^{(1)} \alpha_2^{(2)} \alpha_3^{(3)}}

\noindent von Null verschieden ist, kann in doppelter Weise geometrisch
gedeutet werden, einmal als eine \gesperrt{Abbildung}, indem wir $x_1, x_2, x_3$
und $y_1, y_2, y_3$ als Coordinaten zweier Punkte $x, y$ in demselben
Coordinatensysteme ansehen, so dass jedem Punkte $y$ ein Punkt $x$
entspricht und umgekehrt, sodann auch als \gesperrt{Coordinatentransformation}, wenn wir $x_1, x_2, x_3$ und $y_1, y_2, y_3$ als Coordinaten eines und desselben Punktes, bezogen auf zwei verschiedene
Coordinatensysteme, betrachten. Für die nächsten Betrachtungen
werden wir an der letzteren Interpretation festhalten.
